%------------------------------------------------------------------------------%
\documentclass[fleqn,utf8,aspectratio=169,12pt]{beamer}

\usepackage{integracao_e_series}

\title{Série-$p$}

%------------------------------------------------------------------------------%
\begin{document}

\addtitle

%------------------------------------------------------------------------------%
\section{Série-$p$ e Série Harmônica}

%------------------------------------------------------------------------------%
\begin{frame}{Série-$p$}
  \emph{Série-$p$} ou $p$-série
  \[ \sum_{{\color<2>{magenta}n} =1}^{\infty} \frac{1}{{\color<2>{magenta}n}^p} \;=\;
     \frac{1}{{\color<2>{magenta}1}} +
     \frac{1}{{\color<2>{magenta}2}^p} +
     \frac{1}{{\color<2>{magenta}3}^p} +
     \frac{1}{{\color<2>{magenta}4}^p} +
     \frac{1}{{\color<2>{magenta}5}^p} + \cdots
  \]
  onde $p$ é uma constante real

  \pause
  \vspace{\baselineskip}
  Não confundir com a \emph{Série Geométrica}
  \[
    \sum_{{\color<2>{magenta}n}=1}^{\infty} \alpha r^{\color<2>{magenta}n}
    \;=\; \alpha r^{\color<2>{magenta}0}
    +     \alpha r^{\color<2>{magenta}1}
    +     \alpha r^{\color<2>{magenta}2}
    +     \alpha r^{\color<2>{magenta}3}
    +     \alpha r^{\color<2>{magenta}4}
    + \cdots
  \]
\end{frame}

%------------------------------------------------------------------------------%
\begin{frame}{Série Harmônica}
  \emph{Série Harmônica} é um caso particular da série-$p$ com $p=1$
  \[ \sum_{n=1}^{\infty} \frac{1}{n} \;=\;
     \frac{1}{1} +
     \frac{1}{2} +
     \frac{1}{3} +
     \frac{1}{4} +
     \frac{1}{5} + \cdots
  \]

  \pause
  Essa série passa no teste da divergência, pois,
  \(
    \quad a_n = \frac{1}{\,n\,}\to 0
  \)

  \pause
  \medskip
  mas mesmo assim diverge
\end{frame}

%------------------------------------------------------------------------------%
\begin{proofframe}{Somas Parciais da Série Harmônica}{}
	{\color<3->{gray}
	\[
			S_2 \;=\; 1 + \frac{1}{2} \pause \;=\; \frac{3}{2}
	\]}
	{\color<6->{gray}
	\[
  	\pause
		S_4 \;=\; S_2 + \frac{1}{3} + \frac{1}{4} \pause
	      \;>\; S_2 + \frac{1}{4} + \frac{1}{4} \pause
	      \;=\; \frac{3}{2} + \frac{1}{2}
	      \;=\; \frac{4}{2}
	\]}
	{\color<9->{gray}
	\[
    \pause
    S_8 \;=\; S_4 + \frac{1}{5} + \frac{1}{6} + \frac{1}{7} + \frac{1}{8} \pause
	      \;>\; S_4 + \frac{1}{8} + \frac{1}{8} + \frac{1}{8} + \frac{1}{8} \pause
	      \;>\; \frac{4}{2} + \frac{1}{2}
	      \;=\; \frac{5}{2}
	\]}
	{\color<12->{gray}
    \[
    \pause
    S_{16} \;=\; S_8 + \frac{1}{9}  + \frac{1}{10} + \frac{1}{11}
	                   + \frac{1}{12} + \frac{1}{13} + \frac{1}{14}
	                   + \frac{1}{15} + \frac{1}{16} \pause
	         \;>\; S_8 + \frac{8}{16} \pause
	         \;>\; \frac{5}{2} + \frac{1}{2}
	         \;=\; \frac{6}{2}
	\]}
  \pause
  \[
    S_{2^k} \;\geq\; \frac{2+k}{2}
  \]
\end{proofframe}

%------------------------------------------------------------------------------%
\begin{proofframe}{Somas Parciais da Série Harmônica}
  Considerando os termos $\;n\,=\,2^k$ da sequência de somas parciais

  \pause
  observamos que
  \[
    S_{2^k} \;\geq\; \frac{2+k}{2}
  \]

  \pause
  portanto, $S_n\to\infty$ e a série harmônica diverge
\end{proofframe}

%------------------------------------------------------------------------------%
\section{Aplicando o Teste da Integral para $p\neq 1$}

%------------------------------------------------------------------------------%
\begin{question}
  Utilize o teste da integral para verificar a convergência da série $p$
  com $p\neq 1$
\end{question}

%------------------------------------------------------------------------------%
\begin{resolution}{Teste da Integral}
  Observamos que
  \[
  	a_n = \frac{1}{n^p} = f(n) \qquad \text{para} \qquad f(x) = \frac{1}{x^p}
  \]
  \pause
  Para $x\in[1, \infty)$ a \emph{função} $f(x)$ é
  \pause
  \emph{decrescente},
  \pause
  \emph{contínua}
  \pause
  e \emph{positiva}

  \pause
  \bigskip
  Pelo Teste da Integral
    \[
    \sum_{n=1}^{\infty} \frac{1}{n^p} \quad\text{converge}
    \qquad\Leftrightarrow\qquad
    \int_{1}^{\infty} \frac{1}{x^p}\,dx \quad\text{converge}
  \]
\end{resolution}

%------------------------------------------------------------------------------%
\begin{resolution}{Calculando a Integral}
  \begin{align*}
  	\onslide<+->{    \int_{1}^{\infty} \frac{1}{x^p}\,dx}
    \onslide<+->{& = \lim_{b\to\infty} \int_{1}^{b} x^{-p}\,dx    \qquad p \neq 1                    \\[1.5mm]}
  	\onslide<+->{& = \lim_{b\to\infty} \left[ \;\frac{\,x^{1-p}\,}{\,1-p\,}\evalat{1}{b}\;\, \right] \\[1mm]}
  	\onslide<+->{& = \frac{1}{\,1-p\,} \lim_{b\to\infty} \Big( b^{1-p} - 1 \Big)}
  \end{align*}

  \bigskip
  \onslide<+->{A convergência de \(\quad\lim_{b\to\infty} b^{1-p}\quad\) depende do valor de $p$}
\end{resolution}

%------------------------------------------------------------------------------%
\begin{resolution}{Convergência}
  Se \emph{$p>1$}
  \onslide<2->{\[1-p < 0 \qquad \qquad \lim_{b\to\infty} b^{1-p} = 0\]}

  \vspace{\baselineskip}
  Se \emph{$p<1$}
  \onslide<3->{\[1-p > 0 \qquad \qquad \lim_{b\to\infty} b^{1-p} \quad\text{diverge}\]}
\end{resolution}

%------------------------------------------------------------------------------%
\section{Convergência da Série $p$}
%------------------------------------------------------------------------------%
\begin{frame}{Convergência da Série-$p$}
  A série-$p$
  \[
    \sum_{n=1}^\infty\frac{1}{\;n^p\;}
    = 1
    + \frac{1}{2^p}
    + \frac{1}{3^p}
    + \frac{1}{4^p}
    + \frac{1}{5^p}
    + \frac{1}{6^p}
    + \cdots
  \]
  converge para $p>1$ e diverge caso contrário
\end{frame}

%------------------------------------------------------------------------------%
\section{Lista Mínima}

%------------------------------------------------------------------------------%
\begin{frame}{Lista Mínima}
  Estudar as Seção 6.4 da Apostila

  \vspace{\baselineskip}
  Exercícios:

  \vfill
  Atenção: A prova é baseada no livro, não nas apresentações
\end{frame}

%------------------------------------------------------------------------------%
\end{document}
%------------------------------------------------------------------------------%
